\documentclass[11pt,a4paper]{article}
\usepackage[utf8]{inputenc}
\usepackage[T1]{fontenc}
\usepackage[brazil]{babel}
\usepackage{lmodern}
\usepackage{microtype}
\usepackage{amsmath,amssymb}
\usepackage{geometry}
\usepackage{array}
\usepackage{booktabs}
\usepackage{longtable}
\usepackage{enumitem}
\usepackage{xcolor}
\usepackage{titlesec}
\usepackage{titling}
\usepackage{fancyhdr}
\usepackage{graphicx}
\usepackage{tikz}
\usepackage[most]{tcolorbox}
\usepackage{needspace}
\usepackage[bookmarks=false]{hyperref}

\geometry{
  top=2.4cm,
  bottom=2.6cm,
  left=2.25cm,
  right=2.25cm,
  headheight=22pt
}

\definecolor{astroNavy}{HTML}{14213D}
\definecolor{astroBlue}{HTML}{2F6690}
\definecolor{astroCyan}{HTML}{35A7B8}
\definecolor{astroGold}{HTML}{F2A541}
\definecolor{astroInk}{HTML}{232323}
\definecolor{astroMuted}{HTML}{667085}
\definecolor{astroPaper}{HTML}{F7F9FC}
\definecolor{astroLine}{HTML}{D8E2EC}

\hypersetup{
  colorlinks=true,
  linkcolor=astroBlue,
  urlcolor=astroBlue,
  pdftitle={Astrofisica Moderna: Plano de aulas e fichas de estudo},
  pdfauthor={Material didatico em portugues}
}

\setlength{\parindent}{0pt}
\setlength{\parskip}{0.55em}
\setlist{itemsep=0.15em, topsep=0.25em, leftmargin=1.35em}
\renewcommand{\arraystretch}{1.18}
\color{astroInk}
\graphicspath{{figuras/}}

\pagestyle{fancy}
\fancyhf{}
\lhead{\footnotesize\textcolor{astroMuted}{Astrofisica Moderna}}
\rhead{\footnotesize\textcolor{astroMuted}{\nouppercase{\leftmark}}}
\cfoot{\small\textcolor{astroMuted}{\thepage}}
\renewcommand{\headrulewidth}{0.4pt}
\renewcommand{\headrule}{\hbox to\headwidth{\color{astroLine}\leaders\hrule height \headrulewidth\hfill}}

\titleformat{\section}
  {\Large\bfseries\color{astroNavy}}
  {}
  {0pt}
  {}
  [\vspace{-0.45em}\textcolor{astroGold}{\titlerule[1.1pt]}]
\titlespacing*{\section}{0pt}{1.4em}{0.85em}

\newtcolorbox{noticebox}{
  enhanced,
  breakable,
  colback=astroPaper,
  colframe=astroBlue,
  borderline west={2.4pt}{0pt}{astroGold},
  boxrule=0.5pt,
  arc=2pt,
  left=10pt,
  right=10pt,
  top=8pt,
  bottom=8pt
}

\newtcolorbox{aulahead}[3]{
  enhanced,
  colback=astroNavy,
  colframe=astroNavy,
  arc=2pt,
  boxrule=0pt,
  left=12pt,
  right=12pt,
  top=10pt,
  bottom=10pt,
  overlay={
    \fill[astroGold] (frame.north west) rectangle ([xshift=4pt]frame.south west);
  },
  fontupper=\color{white},
  before skip=1.2em,
  after skip=0.8em,
  title={\large Aula #1},
  fonttitle=\bfseries\color{astroGold},
  attach boxed title to top left={xshift=9pt,yshift=-2pt},
  boxed title style={colback=astroNavy, colframe=astroNavy, boxrule=0pt}
}

\newtcolorbox{lessonblock}[2][]{
  enhanced,
  breakable,
  colback=white,
  colframe=astroLine,
  coltitle=astroNavy,
  colbacktitle=astroPaper,
  boxrule=0.5pt,
  arc=2pt,
  left=10pt,
  right=10pt,
  top=7pt,
  bottom=7pt,
  title={#2},
  fonttitle=\bfseries,
  borderline west={2pt}{0pt}{astroCyan},
  #1
}

\newcommand{\aula}[3]{%
  \Needspace{10\baselineskip}%
  \section*{Aula #1 -- #2}%
  \phantomsection%
  \addcontentsline{toc}{section}{\texorpdfstring{Aula #1 -- #2}{Aula #1 - #2}}%
  \markboth{Aula #1}{Aula #1}%
  \begin{aulahead}{#1}{#2}{#3}%
    \Large\bfseries #2\par\vspace{0.35em}%
    {\small\textcolor{astroLine}{Referência principal no Carroll \& Ostlie: #3}}%
  \end{aulahead}%
  \begin{lessonblock}[borderline west={2pt}{0pt}{astroGold}]{Pergunta inicial}%
  Que observação, escala física ou impasse conceitual torna necessário estudar \textbf{#2}? Antes de ler os objetivos, escreva uma hipótese curta. Ao final da aula, volte a ela e corrija o que mudou.%
  \end{lessonblock}%
}
\newcommand{\objetivos}[1]{\begin{lessonblock}{Ao final, você deve conseguir}#1\end{lessonblock}}
\newcommand{\ideias}[1]{\begin{lessonblock}{Ideias-chave}\begin{itemize}#1\end{itemize}\end{lessonblock}}
\newcommand{\atividade}[1]{\begin{lessonblock}[borderline west={2pt}{0pt}{astroGold}]{Para praticar}#1\end{lessonblock}}
\newcommand{\exercicios}[1]{\begin{lessonblock}[borderline west={2pt}{0pt}{astroBlue}]{Exercícios}\begin{enumerate}#1\end{enumerate}\end{lessonblock}\medskip}
\newcommand{\observaveis}[4]{%
  \begin{lessonblock}[borderline west={2pt}{0pt}{astroBlue}]{Da observação ao modelo}%
  \begin{tabular}{>{\raggedright\arraybackslash}p{0.27\textwidth}>{\raggedright\arraybackslash}p{0.27\textwidth}>{\raggedright\arraybackslash}p{0.34\textwidth}}
  \toprule
  Observável & Grandeza inferida & Hipótese ou calibração necessária \\
  \midrule
  #1 & #2 & #3 \\
  \bottomrule
  \end{tabular}

  \textbf{Pergunta de checagem:} #4
  \end{lessonblock}
}
\newcommand{\aberturabloco}[3]{%
  \begin{noticebox}
  \textbf{Bloco #1.} #2

  \textbf{Como estudar este bloco:} #3
  \end{noticebox}
}
\newcommand{\imagemastro}[4]{%
  \begin{figure}[htbp]
  \centering
  \includegraphics[width=#1\textwidth]{#2}
  \caption{\textbf{#3.} #4}
  \end{figure}
}
\newcommand{\roteirocapitulo}[6]{%
  \newpage
  \begin{lessonblock}[borderline west={2pt}{0pt}{astroGold}]{Estudo dirigido do capítulo}
  \textbf{Tema central:} #1

  \textbf{Mapa conceitual mínimo.} Antes de avançar, escreva uma versão própria dos conceitos a seguir e conecte-os com setas causais, não apenas com setas de associação:
  \begin{itemize}
  #2
  \end{itemize}
  O mapa deve indicar quais grandezas são observáveis diretamente, quais são inferidas por modelo e quais dependem de calibração. Esse cuidado evita uma confusão comum em astrofísica: tratar uma quantidade derivada como se fosse uma medida direta.

  \textbf{Leitura ativa.} Releia o capítulo procurando três frases que expressem relações de causa e efeito. Reescreva cada uma delas no formato ``se ... então ... porque ...''. Esse exercício obriga a explicitar o mecanismo físico, e não apenas o nome do fenômeno.
  \end{lessonblock}

  \begin{lessonblock}{Problema orientado}
  #3

  \textbf{Roteiro de resolução.}
  \begin{enumerate}
  \item Identifique as grandezas dadas e escreva suas unidades.
  \item Desenhe um esquema simples do sistema físico antes de substituir números.
  \item Escolha a relação física apropriada e justifique por que ela se aplica.
  \item Faça a conta mantendo potências de dez separadas dos fatores numéricos.
  \item Interprete o resultado em palavras: ele é grande, pequeno ou plausível?
  \end{enumerate}
  Ao final, escreva uma frase sobre a principal hipótese simplificadora usada. Em astronomia, saber o que foi ignorado é parte da resposta.
  \end{lessonblock}

  \newpage
  \begin{lessonblock}[borderline west={2pt}{0pt}{astroCyan}]{Análise de dados ou gráfico}
  #4

  \textbf{Questões de interpretação.}
  \begin{enumerate}
  \item Qual é a tendência principal que o gráfico ou tabela mostra?
  \item Há algum ponto que parece discrepante? Que explicações físicas ou observacionais poderiam justificá-lo?
  \item Que transformação de escala ajudaria: linear, logarítmica, normalizada ou razão entre grandezas?
  \item Que conclusão seria precipitada sem uma estimativa de incerteza?
  \end{enumerate}
  A resposta deve separar descrição, inferência e especulação. Descrição é o que se vê; inferência é o que se calcula ou modela; especulação é uma hipótese ainda não testada com os dados disponíveis.
  \end{lessonblock}

  \begin{lessonblock}{Erros comuns e antídotos}
  \begin{itemize}
  #5
  \end{itemize}
  Para cada erro, escreva um exemplo concreto e uma pergunta de checagem. A pergunta de checagem deve ser curta o suficiente para ser usada durante uma revisão ou resolução de exercícios.
  \end{lessonblock}

  \newpage
  \begin{lessonblock}[borderline west={2pt}{0pt}{astroBlue}]{Produção escrita}
  #6

  \textbf{Formato sugerido.} Escreva entre 20 e 30 linhas. O primeiro parágrafo deve apresentar o fenômeno; o segundo deve explicar o mecanismo físico; o terceiro deve discutir uma limitação observacional ou uma hipótese do modelo. Evite começar com definições de dicionário. Comece com uma pergunta física ou com uma observação concreta.
  \end{lessonblock}

  \begin{lessonblock}{Exercícios adicionais}
  \begin{enumerate}
  \item Crie uma estimativa de ordem de grandeza relacionada ao capítulo e resolva-a com no máximo três equações.
  \item Escolha uma afirmação qualitativa do texto e traduza-a em uma relação matemática aproximada.
  \item Identifique uma situação em que duas interpretações diferentes poderiam produzir observações parecidas.
  \item Proponha uma observação adicional que ajudaria a distinguir essas interpretações.
  \item Escreva uma pergunta conceitual que ajude a verificar se você entendeu a ideia central.
  \end{enumerate}
  \end{lessonblock}
}

\title{\textbf{Astrofísica Moderna}\\Plano de aulas e fichas de estudo\\48 aulas}
\author{Material didático em português baseado na organização conceitual de\\Bradley W. Carroll e Dale A. Ostlie, \textit{An Introduction to Modern Astrophysics}, 2ª ed.}
\date{2026}

\begin{document}
\pagenumbering{Roman}
\begin{titlepage}
\thispagestyle{empty}
\vspace*{0.8cm}
{\color{astroGold}\rule{\textwidth}{1.2pt}}\par
\vspace{1.0cm}
{\fontsize{32}{36}\selectfont\bfseries\color{astroNavy} Astrofísica Moderna\par}
\vspace{0.35cm}
{\Large\color{astroBlue} Plano de aulas e fichas de estudo\par}
{\large\color{astroMuted} 48 aulas\par}
\vfill
\begin{noticebox}
Este documento reúne as fichas de estudo aula a aula. A apostila principal fica reservada aos capítulos de conteúdo; este arquivo serve como roteiro de acompanhamento, preparação e revisão de cada encontro.
\end{noticebox}
\vfill
{\color{astroMuted}2026\par}
\vspace{0.6cm}
{\color{astroGold}\rule{0.38\textwidth}{1.2pt}}\par
\end{titlepage}
\pagenumbering{roman}

\tableofcontents
\newpage
\pagenumbering{arabic}

\section*{Como usar este plano de aulas}
\addcontentsline{toc}{section}{Como usar este plano de aulas}
Este documento é o roteiro de estudo do semestre. Cada ficha apresenta uma pergunta inicial, objetivos, ideias-chave, uma atividade curta e exercícios. Use a pergunta inicial como gatilho: tente responder antes da aula, volte a ela depois e observe o que mudou.

A leitura conceitual completa está na apostila principal, \texttt{apostila\_astromoderna\_48\_aulas.tex}. Aqui, a função das fichas é organizar o ritmo do curso e indicar o que preparar, revisar e praticar em cada encontro.

\section*{Distribuição das 48 aulas}
\addcontentsline{toc}{section}{Distribuição das 48 aulas}
\begin{longtable}{>{\raggedright\arraybackslash}p{1.2cm}>{\raggedright\arraybackslash}p{9.2cm}>{\raggedright\arraybackslash}p{3.2cm}}
\toprule
Aulas & Tema & Bloco \\
\midrule
1--6 & Ferramentas observacionais: céu, órbitas, luz e telescópios & Fundamentos \\
7--21 & Estrelas: parâmetros, espectros, interiores, Sol, formação e evolução & Astrofísica estelar \\
22--25 & Compactos, relatividade e binárias interagentes & Altas energias \\
26--34 & Sistema Solar, corpos menores e exoplanetas & Astrofísica planetária \\
35--42 & Via Láctea, galáxias, evolução galáctica e estrutura cósmica & Galáxias \\
43--48 & Galáxias ativas, cosmologia, CMB, energia escura e universo primordial & Cosmologia \\
\bottomrule
\end{longtable}

\newpage
\aula{1}{O céu como laboratório físico}{cap. 1}
\objetivos{Reconhecer a astronomia como ciência observacional e conectar fenômenos aparentes do céu a modelos físicos.}
\ideias{
\item A esfera celeste é uma construção geométrica útil, não uma superfície real: coordenadas celestes fixam direção, não posição tridimensional.
\item Medir uma direção é medir um ângulo. A constante $206265$ converte radianos em segundos de arco, e o céu inteiro tem $41\,253$ graus quadrados.
\item Existem quatro sistemas de coordenadas em uso, e a escolha entre eles é operacional: horizontal para saber se dá para observar, equatorial para catalogar, eclíptico para o Sistema Solar, galáctico para a Via Láctea.
\item A lei dos cossenos esférica resolve a conversão entre sistemas e produz a equação da altura, $\sin h=\sin\varphi\sin\delta+\cos\varphi\cos\delta\cos H$, da qual saem culminação, arco diurno e circumpolaridade.
\item As estações vêm da obliquidade de $23^\circ26'$, não da excentricidade: esta produz apenas $7\%$ de variação de fluxo, e com o sinal trocado entre hemisférios.
\item O eixo da Terra precessa a $50{,}3''$ por ano, período de $25\,800$ anos, porque Sol e Lua exercem torque sobre o bojo equatorial. É por isso que todo catálogo precisa declarar sua época.
\item O sistema de referência moderno, o ICRS, é ancorado em quasares, e não no equador terrestre --- e por isso não precisa de correção de precessão.
}

\begin{lessonblock}{Roteiro de 50 minutos}
\begin{tabular}{>{\raggedright\arraybackslash}p{0.13\textwidth}>{\raggedright\arraybackslash}p{0.79\textwidth}}
\toprule
Tempo & Atividade \\
\midrule
0--8 min & Abertura: por que a astronomia é uma ciência de inferência à distância. A cadeia observável $\rightarrow$ grandeza física. \\
8--18 min & Ângulos, pequenos ângulos e ângulo sólido. Calibrar intuição: olho ($1'$), Lua ($0{,}5^\circ$), difração de $8\,\mathrm{m}$ ($0{,}015''$), VLBI ($20\,\mu$as). \\
18--32 min & Os quatro sistemas de coordenadas; dedução da equação da altura pela lei dos cossenos esférica; aplicação imediata à latitude local. \\
32--42 min & Eclíptica, obliquidade e estações, com o argumento quantitativo de $\sin h$ e do arco diurno. Mencionar a equação do tempo. \\
42--50 min & Precessão e nutação: por que catálogos têm época; ICRS e Gaia-CRF; a cadeia de reduções da posição. \\
\bottomrule
\end{tabular}

\textbf{Leitura de apoio:} Capítulo 1 da apostila, seções 1 a 9 (de ``Medir uma direção'' até ``Da posição de catálogo à posição no telescópio'').
\end{lessonblock}
\atividade{Desenhe a esfera celeste local: horizonte, zênite, polos celestes e equador celeste para a latitude de sua cidade.}
\exercicios{
\item Explique por que estrelas diferentes são visíveis em épocas diferentes do ano.
\item Qual é a altura aproximada do polo celeste sul para um observador em Joinville, SC, latitude cerca de $26^\circ$ S?
}

\aula{2}{Coordenadas celestes e tempo astronômico}{cap. 1}
\objetivos{Usar ascensão reta, declinação, altitude e azimute; entender a ligação entre tempo sideral e observação.}
\ideias{
\item O dia solar excede o sideral em $3^{\rm m}56^{\rm s}$ porque a Terra avança $0{,}986^\circ$ em sua órbita a cada dia: o tempo sideral é o relógio que aponta o céu.
\item A paralaxe é geometria pura, $d\,[\mathrm{pc}]=1/p\,['']$, e o parsec é definido exatamente para que essa inversão valha. Desde 2012 a UA é exata por convenção, e o parsec é consequência dela.
\item Aberração ($20{,}5''$) e paralaxe (menos de $1''$) são efeitos distintos, defasados de um quarto de período, e só o segundo depende da distância.
\item Movimento próprio mais distância dão velocidade tangencial, $v_t=4{,}74\,\mu\,d$; com velocidade radial fecha-se o vetor velocidade completo.
\item Inverter a paralaxe é enviesado quando $\sigma_p/p\gtrsim0{,}2$, por efeito de volume --- o viés de Lutz-Kelker.
\item A escala de magnitudes é logarítmica por construção: $m_2-m_1=-2{,}5\log_{10}(F_2/F_1)$, com o $2{,}5$ forçado pela escolha de que cinco magnitudes valem um fator 100.
\item O módulo de distância, $m-M=5\log_{10}d-5+A$, é a ponte entre brilho observado e distância --- e o termo $A$ de extinção não é opcional no plano galáctico.
}

\begin{lessonblock}{Roteiro de 50 minutos}
\begin{tabular}{>{\raggedright\arraybackslash}p{0.13\textwidth}>{\raggedright\arraybackslash}p{0.79\textwidth}}
\toprule
Tempo & Atividade \\
\midrule
0--8 min & Tempo sideral contra tempo solar, deduzindo os $3^{\rm m}56^{\rm s}$; para que serve o tempo sideral na prática observacional. \\
8--20 min & Paralaxe: geometria, definição do parsec, ordem de grandeza dos ângulos envolvidos. Unidades modernas e por que a UA virou exata. \\
20--30 min & Movimento próprio, velocidade tangencial e velocidade espacial; dedução do fator $4{,}74$; viés estatístico ao inverter paralaxes. \\
30--42 min & Fluxo, luminosidade e lei do inverso do quadrado; dedução da escala de Pogson e do módulo de distância. \\
42--50 min & Cores, extinção e avermelhamento: $A_\lambda=1{,}086\,\tau_\lambda$, excesso de cor e $R_V$. Fechamento com a cadeia completa observável $\rightarrow$ grandeza física. \\
\bottomrule
\end{tabular}

\textbf{Leitura de apoio:} Capítulo 1 da apostila, seções 5 e 10 a 16, mais o Laboratório de dados 1.
\end{lessonblock}
\atividade{Escolher uma estrela brilhante e interpretar sua posição em um mapa celeste ou aplicativo planetário.}
\exercicios{
\item Por que a ascensão reta é medida em horas?
\item Uma estrela no equador celeste fica quanto tempo acima do horizonte, aproximadamente, para um observador fora dos polos?
}

\aula{3}{Órbitas elípticas e leis de Kepler}{cap. 2}
\objetivos{Interpretar as leis de Kepler e aplicá-las a planetas, satélites e estrelas binárias.}
\ideias{
\item A elipse possui semieixo maior $a$, excentricidade $e$ e foco ocupado pelo corpo central em órbitas keplerianas.
\item A segunda lei de Kepler expressa conservação do momento angular.
\item A terceira lei relaciona escala orbital e período: $P^2\propto a^3$ para massas fixas.
}
\atividade{Comparar órbitas com $e=0$, $0{,}3$ e $0{,}8$ e discutir onde a velocidade orbital é maior.}
\exercicios{
\item Se $a$ dobra em torno do Sol, por qual fator o período orbital muda?
\item O que aconteceria com a velocidade de um planeta ao passar pelo periélio?
}

\aula{4}{Gravitação newtoniana e teorema do virial}{cap. 2}
\objetivos{Relacionar força gravitacional, energia orbital e sistemas autogravitantes.}
\ideias{
\item A gravitação universal explica queda, órbitas planetárias e dinâmica de aglomerados.
\item Em órbitas ligadas, energia total é negativa.
\item Para muitos sistemas em equilíbrio, o teorema do virial dá $2K+U=0$.
}
\atividade{Estimar a massa do Sol a partir de $P=1$ ano e $a=1$ UA.}
\exercicios{
\item Derive a forma proporcional $M\propto a^3/P^2$ para uma órbita circular.
\item Por que aquecer um sistema autogravitante pode fazê-lo expandir?
}

\aula{5}{Luz, magnitudes e corpo negro}{cap. 3}
\objetivos{Usar brilho aparente, fluxo, magnitude e radiação de corpo negro para extrair informação física.}
\ideias{
\item Fluxo diminui com a distância como $F=L/(4\pi d^2)$.
\item A escala de magnitudes é logarítmica: diferença de 5 magnitudes equivale a fator 100 em fluxo.
\item A lei de Wien, $\lambda_{\max}T\simeq2{,}9\times10^{-3}\,\mathrm{m\,K}$, liga cor e temperatura.
}
\atividade{Comparar espectros de corpos negros de 3000 K, 5800 K e 10000 K.}
\exercicios{
\item Uma estrela é 100 vezes mais brilhante que outra. Qual é a diferença de magnitude?
\item Estime o comprimento de onda de máximo de emissão do Sol usando $T\approx5800$ K.
}

\aula{6}{Fótons, espectros e telescópios}{caps. 5 e 6}
\objetivos{Compreender como linhas espectrais e instrumentos transformam luz em diagnóstico astrofísico.}
\ideias{
\item Energia de fóton: $E=h\nu=hc/\lambda$.
\item Linhas de absorção e emissão revelam composição, temperatura, densidade e velocidade.
\item A resolução angular de um telescópio melhora com abertura maior e comprimento de onda menor.
}
\atividade{Identificar, em um espectro simples, contínuo, linhas de absorção e linhas de emissão.}
\exercicios{
\item Um fóton azul tem energia maior ou menor que um fóton vermelho? Justifique.
\item Por que rádio telescópios costumam ter antenas grandes?
}

\aula{7}{Sistemas binários e massas estelares}{cap. 7}
\objetivos{Mostrar por que estrelas binárias são fundamentais para medir massas.}
\ideias{
\item Binárias visuais, espectroscópicas e eclipsantes oferecem informações complementares.
\item Curvas de velocidade radial revelam movimento ao longo da linha de visada.
\item Massa é o parâmetro mais importante para a evolução estelar.
}
\atividade{Interpretar uma curva de luz de binária eclipsante e localizar eclipse primário e secundário.}
\exercicios{
\item Por que estrelas isoladas têm massas mais difíceis de medir?
\item Que informação adicional uma binária eclipsante fornece em relação a uma binária apenas espectroscópica?
}

\aula{8}{Classificação espectral e diagrama HR}{cap. 8}
\objetivos{Ler o diagrama Hertzsprung-Russell e relacionar classe espectral, temperatura, luminosidade e raio.}
\ideias{
\item A sequência OBAFGKM organiza estrelas por temperatura decrescente.
\item No diagrama HR, a sequência principal reflete fusão estável de hidrogênio no núcleo.
\item Gigantes e anãs brancas ocupam regiões distintas por terem raios muito diferentes.
}
\atividade{Posicionar Sol, Sirius, Betelgeuse e uma anã branca típica em um diagrama HR esquemático.}
\exercicios{
\item Duas estrelas têm a mesma temperatura, mas uma é 100 vezes mais luminosa. Qual deve ter raio maior?
\item Por que a cor de uma estrela é um indicador imperfeito, mas útil, de temperatura?
}

\aula{9}{Atmosferas estelares e opacidade}{cap. 9}
\objetivos{Entender por que a radiação que observamos vem de camadas externas e como a opacidade molda espectros.}
\ideias{
\item A fotosfera é aproximadamente a região de profundidade óptica $\tau\sim1$.
\item Opacidade mede quão eficientemente a matéria bloqueia ou absorve radiação.
\item Perfis de linha dependem de temperatura, pressão, rotação e campos magnéticos.
}
\atividade{Discutir por que não vemos diretamente o núcleo do Sol, embora a energia venha de lá.}
\exercicios{
\item Defina, em palavras, profundidade óptica.
\item Cite dois mecanismos que podem alargar uma linha espectral.
}

\aula{10}{Equilíbrio hidrostático}{cap. 10}
\objetivos{Derivar a ideia física de equilíbrio entre gravidade e gradiente de pressão no interior estelar.}
\ideias{
\item Estrelas são esferas de plasma sustentadas por pressão interna.
\item Em equilíbrio, $dP/dr=-Gm(r)\rho/r^2$.
\item A pressão precisa crescer para dentro porque o peso das camadas superiores aumenta.
}
\atividade{Fazer uma analogia quantitativa com pressão no fundo de uma coluna de fluido.}
\exercicios{
\item O que ocorreria se a pressão interna de uma estrela diminuísse subitamente?
\item Por que o centro de uma estrela tem temperatura maior que a superfície?
}

\aula{11}{Fontes de energia das estrelas}{cap. 10}
\objetivos{Comparar contração gravitacional, reações nucleares e escalas de tempo estelares.}
\ideias{
\item A fusão de hidrogênio em hélio é a fonte dominante na sequência principal.
\item Cadeia próton-próton domina em estrelas como o Sol; ciclo CNO cresce em importância em estrelas mais massivas.
\item Pequenas diferenças de massa nuclear aparecem como energia via $E=mc^2$.
}
\atividade{Estimar a fração de massa convertida em energia na fusão H$\rightarrow$He e discutir a longevidade solar.}
\exercicios{
\item Por que reações nucleares exigem temperaturas tão altas?
\item Explique por que estrelas mais massivas vivem menos, apesar de terem mais combustível.
}

\aula{12}{Transporte de energia no interior estelar}{cap. 10}
\objetivos{Distinguir transporte radiativo, convectivo e condutivo em estrelas.}
\ideias{
\item Transporte radiativo ocorre por absorções e reemissões sucessivas.
\item Convecção aparece quando o gradiente de temperatura torna o fluido instável.
\item Condutividade é crucial em matéria degenerada, como anãs brancas.
}
\atividade{Comparar água aquecida numa panela com uma camada estelar convectiva.}
\exercicios{
\item Por que fótons podem levar muito tempo para escapar do interior solar?
\item Que condição qualitativa favorece convecção?
}

\aula{13}{Modelos estelares e sequência principal}{cap. 10}
\objetivos{Entender o conjunto mínimo de equações de estrutura estelar e a relação massa-luminosidade.}
\ideias{
\item Modelos estelares combinam massa, pressão, luminosidade e temperatura como funções do raio.
\item Para estrelas de sequência principal, luminosidade cresce rapidamente com a massa.
\item A posição no HR é resultado da estrutura interna e da taxa de geração de energia.
}
\atividade{Esboçar como muda uma estrela de sequência principal ao aumentar a massa.}
\exercicios{
\item Se $L\propto M^{3{,}5}$, uma estrela de $2M_\odot$ é quantas vezes mais luminosa que o Sol?
\item Por que a sequência principal não é uma linha de idade única?
}

\aula{14}{O Sol como estrela padrão}{cap. 11}
\objetivos{Usar o Sol como laboratório para estrutura, atmosfera e atividade magnética estelar.}
\ideias{
\item A fotosfera produz o contínuo visível; cromosfera e coroa mostram plasma magnetizado.
\item Manchas solares são regiões mais frias associadas a campos magnéticos intensos.
\item Neutrinos e heliossismologia sondam o interior solar.
}
\atividade{Analisar uma imagem solar em luz branca ou H-alfa e identificar estruturas.}
\exercicios{
\item Por que a coroa solar pode ser mais quente que a fotosfera?
\item O que a detecção de neutrinos solares confirma sobre o núcleo?
}

\aula{15}{Meio interestelar: gás e poeira}{cap. 12}
\objetivos{Descrever fases do meio interestelar e seus efeitos sobre observações.}
\ideias{
\item O meio interestelar contém gás ionizado, atômico, molecular e poeira.
\item Poeira causa extinção e avermelhamento da luz estelar.
\item Linhas como 21 cm do hidrogênio neutro mapeiam estrutura galáctica.
}
\atividade{Comparar a cor observada de uma estrela sem poeira e com avermelhamento interestelar.}
\exercicios{
\item Diferencie extinção e avermelhamento.
\item Por que nuvens moleculares são locais naturais de formação estelar?
}

\aula{16}{Formação estelar e protoestrelas}{cap. 12}
\objetivos{Explicar colapso gravitacional, critério de Jeans e evolução pré-sequência principal.}
\ideias{
\item Uma nuvem colapsa quando a gravidade vence pressão, turbulência e campos magnéticos.
\item A massa de Jeans expressa uma escala crítica de instabilidade.
\item Discos circumestelares surgem por conservação de momento angular.
}
\atividade{Estimar qualitativamente por que uma nuvem ao contrair tende a girar mais rápido.}
\exercicios{
\item O que impede uma nuvem qualquer de colapsar imediatamente?
\item Por que estrelas jovens costumam estar associadas a nebulosas?
}

\aula{17}{Evolução na sequência principal}{cap. 13}
\objetivos{Relacionar massa inicial, composição e tempo de vida durante a queima de hidrogênio.}
\ideias{
\item A sequência principal é uma fase longa de equilíbrio térmico e hidrostático.
\item A composição central muda com a fusão, alterando pressão, temperatura e luminosidade.
\item Massa inicial define grande parte do destino evolutivo.
}
\atividade{Comparar trajetórias evolutivas de estrelas de baixa e alta massa no HR.}
\exercicios{
\item Por que aglomerados estelares são úteis para testar evolução estelar?
\item O que significa ponto de saída da sequência principal em um aglomerado?
}

\aula{18}{Gigantes, queima de hélio e estágios tardios}{cap. 13}
\objetivos{Descrever a evolução pós-sequência principal de estrelas de baixa e massa intermediária.}
\ideias{
\item Ao esgotar hidrogênio central, a estrela reajusta sua estrutura e pode virar gigante vermelha.
\item Queima de hélio produz carbono e oxigênio.
\item Perda de massa e nebulosas planetárias antecedem anãs brancas em muitos casos.
}
\atividade{Interpretar o HR de um aglomerado com ramo de gigantes.}
\exercicios{
\item Por que uma estrela expande quando o núcleo contrai?
\item Uma nebulosa planetária tem relação direta com planetas? Explique.
}

\aula{19}{Pulsação estelar e sismologia}{cap. 14}
\objetivos{Entender pulsação como ferramenta de distância e de estrutura interna.}
\ideias{
\item Cefeidas obedecem uma relação período-luminosidade essencial para a escala de distâncias.
\item Pulsação resulta de mecanismos de armazenamento e liberação de energia em camadas ionizadas.
\item Heliossismologia e asterossismologia usam oscilações para sondar interiores.
}
\atividade{Usar uma relação período-luminosidade simplificada para estimar distância de uma Cefeida.}
\exercicios{
\item Por que Cefeidas são chamadas de velas padrão?
\item Qual é a vantagem de medir frequências de oscilação em uma estrela?
}

\aula{20}{Estrelas massivas e supernovas}{cap. 15}
\objetivos{Explicar evolução de estrelas massivas, colapso de núcleo e produção de elementos.}
\ideias{
\item Estrelas massivas podem queimar sucessivamente elementos até formar núcleos ricos em ferro.
\item Ferro não libera energia por fusão em condições estelares comuns.
\item Supernovas de colapso de núcleo enriquecem o meio interestelar e deixam remanescentes compactos.
}
\atividade{Montar o ``modelo de cebola'' de uma estrela massiva evoluída.}
\exercicios{
\item Por que o ferro marca uma fronteira energética na fusão nuclear?
\item Cite dois efeitos astrofísicos de uma supernova na galáxia.
}

\aula{21}{Anãs brancas e matéria degenerada}{cap. 16}
\objetivos{Introduzir pressão de degenerescência e limite de Chandrasekhar.}
\ideias{
\item Anãs brancas são remanescentes sustentados por pressão de degenerescência dos elétrons.
\item A massa máxima típica é próxima de $1{,}4M_\odot$.
\item Elas esfriam lentamente, funcionando como fósseis térmicos de populações estelares.
}
\atividade{Comparar qualitativamente pressão térmica e pressão de degenerescência.}
\exercicios{
\item Por que uma anã branca não precisa de fusão ativa para se sustentar?
\item O que pode ocorrer se uma anã branca em binária ganhar massa demais?
}

\aula{22}{Estrelas de nêutrons e pulsares}{cap. 16}
\objetivos{Descrever remanescentes compactos, conservação de momento angular e emissão pulsada.}
\ideias{
\item Estrelas de nêutrons têm massas solares comprimidas em raios de ordem dezenas de quilômetros.
\item Pulsares são estrelas de nêutrons magnetizadas em rotação.
\item Campos magnéticos intensos e rotação rápida tornam esses objetos laboratórios extremos.
}
\atividade{Estimar o aumento de rotação quando um núcleo estelar colapsa conservando momento angular.}
\exercicios{
\item Por que pulsares parecem piscar?
\item Qual propriedade torna estrelas de nêutrons muito mais densas que anãs brancas?
}

\aula{23}{Relatividade geral e buracos negros}{cap. 17}
\objetivos{Apresentar gravidade como geometria e definir horizonte de eventos e raio de Schwarzschild.}
\ideias{
\item Relatividade geral descreve gravidade como curvatura do espaço-tempo.
\item O raio de Schwarzschild é $R_s=2GM/c^2$.
\item Buracos negros são detectados por seus efeitos gravitacionais e por acreção, não por luz emitida do horizonte.
}
\atividade{Calcular $R_s$ para uma massa solar e comparar com o raio do Sol.}
\exercicios{
\item O que é horizonte de eventos?
\item Por que um buraco negro em binária pode ser brilhante em raios X?
}

\aula{24}{Sistemas binários próximos e acreção}{cap. 18}
\objetivos{Entender lóbulos de Roche, transferência de massa e discos de acreção.}
\ideias{
\item Em binárias próximas, uma estrela pode preencher seu lóbulo de Roche.
\item A matéria transferida forma disco quando possui momento angular suficiente.
\item Novas, supernovas Ia e binárias de raios X dependem de acreção sobre objetos compactos.
}
\atividade{Desenhar equipotenciais em uma binária e localizar o ponto de Lagrange interno.}
\exercicios{
\item Por que discos de acreção aquecem?
\item Qual é a diferença geral entre nova e supernova Ia?
}

\aula{25}{Síntese da astrofísica estelar}{caps. 7--18}
\objetivos{Integrar nascimento, vida e morte das estrelas em uma narrativa física coerente.}
\ideias{
\item Massa inicial é a variável que mais organiza a evolução estelar.
\item Observações de espectros, luminosidades e variabilidade testam modelos.
\item Remanescentes compactos conectam evolução estelar, relatividade e altas energias.
}
\atividade{Construir um mapa conceitual: nuvem molecular $\rightarrow$ estrela $\rightarrow$ remanescente.}
\exercicios{
\item Compare os destinos prováveis de estrelas de $1M_\odot$, $8M_\odot$ e $30M_\odot$.
\item Que observações permitiriam estimar idade de um aglomerado?
}

\aula{26}{Processos físicos no Sistema Solar}{cap. 19}
\objetivos{Aplicar gravitação, equilíbrio térmico, atmosferas e marés ao Sistema Solar.}
\ideias{
\item Planetas refletem e emitem radiação; equilíbrio térmico depende de distância, albedo e atmosfera.
\item Marés surgem de diferenças espaciais na força gravitacional.
\item Atmosferas planetárias dependem de gravidade, temperatura, composição e história.
}
\atividade{Comparar temperatura de equilíbrio de Mercúrio, Terra e Marte ignorando atmosfera.}
\exercicios{
\item Por que a Lua mostra quase sempre a mesma face para a Terra?
\item Como o albedo altera a temperatura de equilíbrio?
}

\aula{27}{Planetas terrestres}{cap. 20}
\objetivos{Comparar Mercúrio, Vênus, Terra, Lua e Marte como corpos rochosos.}
\ideias{
\item Planetas terrestres têm superfícies sólidas e histórias geológicas diversas.
\item Efeito estufa explica a temperatura elevada de Vênus.
\item Marte registra evidências de água passada e evolução atmosférica importante.
}
\atividade{Montar uma tabela comparativa de atmosfera, campo magnético, água e atividade geológica.}
\exercicios{
\item Por que Vênus é mais quente que Mercúrio na superfície?
\item Cite uma evidência de água líquida antiga em Marte.
}

\aula{28}{Gigantes gasosos, luas e anéis}{cap. 21}
\objetivos{Descrever estrutura de planetas gigantes e diversidade de satélites.}
\ideias{
\item Júpiter e Saturno são ricos em hidrogênio e hélio; Urano e Netuno têm mais compostos voláteis pesados.
\item Luas gigantes podem ter oceanos internos, vulcanismo ou atmosferas.
\item Anéis são sistemas de partículas moldados por gravidade, colisões e ressonâncias.
}
\atividade{Comparar Io, Europa, Titã e Encélado como mundos ativos.}
\exercicios{
\item Por que Europa é alvo de interesse astrobiológico?
\item O que mantém anéis planetários finos e estruturados?
}

\aula{29}{Corpos menores: asteroides, cometas e cinturão de Kuiper}{cap. 22}
\objetivos{Entender corpos menores como registros da formação do Sistema Solar.}
\ideias{
\item Asteroides concentram-se principalmente entre Marte e Júpiter, mas há muitas populações dinâmicas.
\item Cometas preservam material volátil e desenvolvem coma e cauda ao se aproximarem do Sol.
\item Objetos transnetunianos ajudam a reconstruir a história externa do Sistema Solar.
}
\atividade{Interpretar a diferença entre cauda de poeira e cauda iônica de um cometa.}
\exercicios{
\item Por que meteoritos são amostras importantes?
\item Qual é a diferença entre meteoroide, meteoro e meteorito?
}

\aula{30}{Formação do Sistema Solar}{cap. 23}
\objetivos{Apresentar a hipótese nebular, acreção planetesimal e diferenciação planetária.}
\ideias{
\item O Sistema Solar nasceu de um disco protoplanetário ao redor do jovem Sol.
\item A linha de gelo influenciou a composição dos planetas.
\item Migração planetária e impactos gigantes ajudaram a moldar a arquitetura atual.
}
\atividade{Desenhar um disco com gradiente de temperatura e localizar regiões de planetas rochosos e gigantes.}
\exercicios{
\item Por que planetas gigantes se formaram preferencialmente longe do Sol?
\item Como a conservação de momento angular explica a formação de um disco?
}

\aula{31}{Exoplanetas: métodos de detecção}{caps. 7 e 23}
\objetivos{Explicar velocidade radial, trânsito, imageamento direto, microlente e astrometria.}
\ideias{
\item Velocidade radial mede o balanço gravitacional estrela-planeta.
\item Trânsitos medem queda de fluxo e fornecem raio relativo.
\item Métodos diferentes favorecem tipos diferentes de planetas e órbitas.
}
\atividade{Calcular a profundidade de trânsito aproximada de um planeta do tamanho de Júpiter diante de uma estrela solar.}
\exercicios{
\item Por que planetas grandes e próximos são mais fáceis de detectar por trânsito?
\item Que combinação de medidas permite obter densidade média de um exoplaneta?
}

\aula{32}{Arquiteturas planetárias e habitabilidade}{cap. 23}
\objetivos{Comparar sistemas exoplanetários observados com o Sistema Solar e discutir habitabilidade com cautela.}
\ideias{
\item Sistemas exoplanetários mostram grande diversidade: Júpiteres quentes, super-Terras e cadeias ressonantes.
\item Zona habitável é condição orbital, não garantia de vida.
\item Atmosferas, atividade estelar e história geológica são essenciais para avaliar habitabilidade.
}
\atividade{Debater critérios mínimos para priorizar alvos de busca por atmosferas.}
\exercicios{
\item Por que uma zona habitável depende da luminosidade da estrela?
\item Cite duas razões pelas quais ``habitável'' não significa ``habitado''.
}

\aula{33}{Escalas, populações e inventário planetário}{caps. 19--23}
\objetivos{Consolidar o Sistema Solar e exoplanetas por escalas de massa, distância e composição.}
\ideias{
\item A classificação de objetos planetários envolve composição, dinâmica e história de formação.
\item O Sistema Solar é uma referência, mas não representa toda a diversidade planetária.
\item Pequenos corpos e discos debris conectam sistemas jovens e maduros.
}
\atividade{Construir uma escala logarítmica de distâncias: raio terrestre, UA, cinturão de Kuiper, parsec.}
\exercicios{
\item Por que escalas logarítmicas são úteis em astronomia?
\item Escolha um corpo menor e explique que informação ele traz sobre formação planetária.
}

\aula{34}{Revisão aplicada: do céu local aos sistemas planetários}{caps. 1--23}
\objetivos{Preparar avaliação integrando observação, luz, estrelas e planetas.}
\ideias{
\item A astrofísica combina medidas indiretas: posição, fluxo, espectro, tempo e polarização.
\item Modelos simples permitem estimativas de ordem de grandeza poderosas.
\item A interpretação depende de hipóteses explícitas e incertezas.
}
\atividade{Resolva um problema aberto: estime propriedades de uma estrela com planeta em trânsito a partir de uma curva de luz simplificada.}
\exercicios{
\item Liste cinco grandezas observáveis e cinco grandezas inferidas em astrofísica.
\item Explique por que uma mesma observação pode admitir mais de uma interpretação.
}

\aula{35}{A Via Láctea: estrutura geral}{cap. 24}
\objetivos{Descrever componentes da Galáxia: disco, bojo, halo, braços espirais e matéria escura.}
\ideias{
\item A Via Láctea é uma galáxia espiral barrada.
\item Poeira dificulta observação no visível, mas rádio e infravermelho revelam estrutura.
\item Populações estelares diferentes indicam histórias de formação distintas.
}
\atividade{Desenhar a Via Láctea vista de lado e de cima, com localização aproximada do Sol.}
\exercicios{
\item Por que é difícil mapear a Via Láctea estando dentro dela?
\item Diferencie disco fino, bojo e halo.
}

\aula{36}{Cinemática galáctica e matéria escura}{cap. 24}
\objetivos{Usar curvas de rotação para inferir distribuição de massa.}
\ideias{
\item Velocidades orbitais no disco revelam a massa interna à órbita.
\item Curvas de rotação aproximadamente planas indicam massa não luminosa extensa.
\item Matéria escura é inferida gravitacionalmente em várias escalas.
}
\atividade{Comparar curva esperada para massa concentrada com curva observada plana.}
\exercicios{
\item Se a massa estivesse concentrada no centro, como a velocidade orbital mudaria com o raio?
\item Cite outra evidência de matéria escura além de curvas de rotação.
}

\aula{37}{O centro galáctico}{cap. 24}
\objetivos{Apresentar evidências para o buraco negro supermassivo central da Via Láctea.}
\ideias{
\item Órbitas de estrelas próximas ao centro permitem medir uma massa compacta.
\item Sgr A* é a fonte associada ao buraco negro central.
\item O centro galáctico combina densidade estelar alta, gás, poeira e campos magnéticos.
}
\atividade{Interpretar um gráfico de órbita de uma estrela ao redor de Sgr A*.}
\exercicios{
\item Por que órbitas estelares são evidência forte de massa compacta?
\item Por que observar o centro galáctico no infravermelho é vantajoso?
}

\aula{38}{Tipos de galáxias e sequência de Hubble}{cap. 25}
\objetivos{Classificar galáxias e reconhecer limitações de classificações morfológicas.}
\ideias{
\item Galáxias elípticas, espirais e irregulares diferem em forma, gás e formação estelar.
\item Morfologia não é necessariamente sequência evolutiva simples.
\item Cor integrada ajuda a distinguir populações jovens e antigas.
}
\atividade{Classificar imagens de galáxias e justificar critérios usados.}
\exercicios{
\item Por que galáxias espirais tendem a ter regiões azuis nos braços?
\item O que uma galáxia avermelhada e pobre em gás sugere sobre formação estelar atual?
}

\aula{39}{Estrutura espiral e formação estelar em galáxias}{cap. 25}
\objetivos{Explicar braços espirais como padrões dinâmicos e locais de formação estelar.}
\ideias{
\item Braços espirais podem ser ondas de densidade ou estruturas transientes.
\item Gás comprimido favorece formação de estrelas massivas, brilhantes e azuis.
\item Estrelas velhas e jovens traçam componentes diferentes da galáxia.
}
\atividade{Usar uma analogia de tráfego para discutir onda de densidade sem transportar sempre as mesmas estrelas.}
\exercicios{
\item Por que braços espirais são destacados por estrelas jovens?
\item Qual observação ajudaria a mapear gás molecular em uma galáxia?
}

\aula{40}{Galáxias elípticas, interações e fusões}{caps. 25 e 26}
\objetivos{Relacionar morfologia, dinâmica e evolução galáctica por interações.}
\ideias{
\item Galáxias elípticas costumam ter pouco gás frio e populações estelares antigas.
\item Encontros gravitacionais geram caudas de maré, pontes e surtos de formação estelar.
\item Fusões contribuem para crescimento de galáxias e buracos negros centrais.
}
\atividade{Analisar imagens de galáxias em interação e identificar marcas dinâmicas.}
\exercicios{
\item Por que colisões de galáxias raramente significam colisões diretas entre estrelas?
\item Como uma fusão pode intensificar formação estelar?
}

\aula{41}{Escada de distâncias extragalácticas}{cap. 27}
\objetivos{Compreender paralaxe, Cefeidas, supernovas Ia e redshift como etapas de distância.}
\ideias{
\item Cada indicador de distância calibra o próximo em uma cadeia.
\item Cefeidas e supernovas Ia são importantes por relacionarem luminosidade observável e intrínseca.
\item Incertezas sistemáticas se propagam pela escada de distâncias.
}
\atividade{Montar a sequência: paralaxe $\rightarrow$ Cefeidas $\rightarrow$ supernovas Ia $\rightarrow$ Hubble.}
\exercicios{
\item Por que paralaxe não basta para medir distâncias cosmológicas?
\item O que torna supernovas Ia úteis como indicadores de distância?
}

\aula{42}{Expansão do universo e grandes estruturas}{cap. 27}
\objetivos{Introduzir redshift cosmológico, lei de Hubble-Lemaître e aglomerados de galáxias.}
\ideias{
\item Redshift cosmológico expressa expansão do espaço em grandes escalas.
\item A lei $v=H_0d$ é uma relação estatística para galáxias distantes no fluxo cósmico.
\item Galáxias formam grupos, aglomerados, filamentos e vazios.
}
\atividade{Fazer um modelo com pontos em uma bexiga ou malha elástica para visualizar expansão relativa.}
\exercicios{
\item Por que a expansão não deve ser imaginada como galáxias fugindo de um centro comum?
\item Se $H_0=70\,\mathrm{km\,s^{-1}\,Mpc^{-1}}$, estime a velocidade recessional de uma galáxia a 100 Mpc.
}

\aula{43}{Galáxias ativas e quasares}{cap. 28}
\objetivos{Explicar núcleos ativos como acreção em buracos negros supermassivos.}
\ideias{
\item AGNs podem emitir enorme luminosidade a partir de regiões compactas.
\item O modelo unificado envolve disco de acreção, toro de poeira, jatos e orientação.
\item Quasares permitem estudar o universo distante e o meio intergaláctico.
}
\atividade{Desenhar o modelo unificado de AGN e prever como muda a aparência com a orientação.}
\exercicios{
\item Por que variabilidade rápida indica região emissora pequena?
\item Como acreção pode converter energia gravitacional em radiação?
}

\aula{44}{Cosmologia newtoniana e parâmetros cósmicos}{cap. 29}
\objetivos{Apresentar densidade crítica, curvatura efetiva e componentes do universo.}
\ideias{
\item A densidade crítica separa comportamentos dinâmicos em modelos simples.
\item O universo atual contém matéria bariônica, matéria escura, radiação e energia escura.
\item Parâmetros cosmológicos são medidos combinando observações independentes.
}
\atividade{Construir um gráfico de pizza aproximado dos componentes energéticos do universo atual.}
\exercicios{
\item Por que matéria bariônica não pode ser toda a matéria escura?
\item O que significa dizer que o universo é aproximadamente plano?
}

\aula{45}{Radiação cósmica de fundo}{cap. 29}
\objetivos{Entender a CMB como relíquia térmica do universo jovem.}
\ideias{
\item A CMB tem espectro quase perfeito de corpo negro com temperatura próxima de 2,7 K.
\item Ela foi emitida quando o universo se tornou transparente à radiação.
\item Pequenas anisotropias registram sementes das estruturas cósmicas.
}
\atividade{Comparar o espectro de corpo negro da CMB com o de uma estrela.}
\exercicios{
\item Por que a CMB aparece hoje em micro-ondas?
\item Que informação física está codificada nas anisotropias da CMB?
}

\aula{46}{Cosmologia relativística e energia escura}{cap. 29}
\objetivos{Discutir expansão acelerada, supernovas Ia e interpretação de energia escura.}
\ideias{
\item Em cosmologia moderna, o fator de escala descreve a expansão do universo.
\item Supernovas Ia distantes indicaram expansão acelerada.
\item Energia escura é uma descrição fenomenológica do componente associado à aceleração cósmica.
}
\atividade{Comparar, qualitativamente, universos com desaceleração, expansão constante e aceleração.}
\exercicios{
\item Por que supernovas Ia testam a história de expansão?
\item Qual é a diferença entre matéria escura e energia escura?
}

\aula{47}{Universo primordial e inflação}{cap. 30}
\objetivos{Apresentar nucleossíntese primordial, inflação e origem das estruturas.}
\ideias{
\item No universo quente e denso, partículas, radiação e núcleos leves estiveram em interação intensa.
\item Nucleossíntese primordial produziu principalmente hidrogênio, hélio e traços de lítio.
\item Inflação é proposta para explicar homogeneidade, planura e origem de flutuações.
}
\atividade{Construir uma linha do tempo: inflação, nucleossíntese, recombinação, primeiras estrelas, galáxias.}
\exercicios{
\item Por que elementos pesados não foram produzidos em grande quantidade no Big Bang?
\item Que problema a inflação ajuda a resolver?
}

\aula{48}{Síntese final: uma narrativa física do universo}{todo o livro}
\objetivos{Integrar escalas e métodos da astrofísica moderna em uma visão de conjunto.}
\ideias{
\item A mesma física aparece em escalas muito diferentes: gravidade, radiação, termodinâmica, plasma e partículas.
\item A astrofísica é uma ciência de inferência: observamos luz e movimento para reconstruir propriedades distantes.
\item O universo observável tem história: expansão, formação de estruturas, nascimento e morte de estrelas, enriquecimento químico e formação planetária.
}
\atividade{Escolha um objeto astronômico e responda: como medir distância, massa, composição, idade aproximada e contexto evolutivo?}
\exercicios{
\item Escolha três equações da apostila e explique o papel físico de cada uma.
\item Escreva uma síntese de uma página conectando estrelas, galáxias e cosmologia.
}



\newpage
\section*{Como estudar ao longo do semestre}
\addcontentsline{toc}{section}{Como estudar ao longo do semestre}
\begin{itemize}
\item \textbf{Listas curtas:} resolva problemas de estimativa, interpretação de gráficos e leitura de espectros sem esperar dominar todos os detalhes algébricos.
\item \textbf{Diário observacional:} registre céu, fases da Lua, planetas visíveis, mapas ou simulações. O objetivo é conectar coordenadas e tempo astronômico com experiência concreta.
\item \textbf{Sínteses de uma página:} ao terminar um bloco, escreva uma explicação curta conectando observação, modelo físico e principal equação.
\item \textbf{Revisão 1:} fundamentos, luz, estrelas e Sistema Solar.
\item \textbf{Revisão 2:} galáxias, cosmologia e síntese conceitual.
\item \textbf{Projeto final:} escolha um sistema astrofísico e estude-o com dados públicos simples, sempre separando observação direta, cálculo e interpretação.
\end{itemize}

\section*{Projetos de estudo possíveis}
\addcontentsline{toc}{section}{Projetos de estudo possíveis}
\begin{enumerate}
\item Estimar temperatura de estrelas por cor usando fotometria pública.
\item Construir um diagrama HR de um aglomerado com dados do Gaia.
\item Modelar uma curva de trânsito exoplanetário simplificada.
\item Comparar curvas de rotação e inferir matéria escura em galáxias.
\item Reproduzir uma versão simples da lei de Hubble-Lemaître com catálogo de galáxias.
\item Fazer uma linha do tempo quantitativa do universo, do Big Bang ao Sistema Solar.
\end{enumerate}
\section*{Bibliografia}
\addcontentsline{toc}{section}{Bibliografia}
\begin{itemize}
\item Carroll, B. W.; Ostlie, D. A. \textit{An Introduction to Modern Astrophysics}. 2ª ed. Cambridge University Press, 2017.
\item Karttunen, H. et al. \textit{Fundamental Astronomy}. Springer.
\item Ryden, B.; Peterson, B. M. \textit{Foundations of Astrophysics}. Cambridge University Press.
\item NASA/IPAC Extragalactic Database, SIMBAD, ESA Gaia Archive e NASA Exoplanet Archive para consultas orientadas de dados.
\end{itemize}

\end{document}
